\section{Correcting the earlier version}
\label{sec:erratum}

This study began as an evaluation of the ONRS composite on FinQA using 4-bit models on a consumer GPU.
Re-running it with more items and a test split exposed three problems, which we report because the earlier numbers are public.
\begin{enumerate}[leftmargin=*,itemsep=2pt,topsep=2pt]
\item \textbf{A scoring flaw.} The evidence-support score counted every number in the response, including the one in the model's own ``Confidence: 0.9'' line, which no evidence can support. Every response carried one extra unsupported number, lowering all support scores and pushing some under the hallucination threshold on that number alone. After removing the confidence statement, the hallucination rates reported earlier ($0.08$ in the baseline, $0.10\rightarrow0.20$ under noise) are $0.00$ in every condition.
\item \textbf{A manipulation that did not work.} The noise experiment added up to four sentences about unrelated figures. At 150 items per split the paired per-item change in accuracy between $0$ and $40\%$ noise was $+0.007$ (95\% interval $-0.05$ to $+0.07$) on dev and $-0.027$ ($-0.08$ to $+0.02$) on test. The injected text reached the model but did not hurt it, so a metric could not be faulted for failing to track it. Corrupting the numbers of the evidence (Section~\ref{sec:setup}) is the replacement.
\item \textbf{Unstable summary statistics.} The original conclusion that the composite fails to track noise and that no weighting of its four components rescues it came from a correlation over five noise-level means and a sweep of 1{,}771 weightings. On the same 3B model the correlation is $+0.09$ on dev and $-0.64$ on test, each with an interval of about $\pm0.8$, and the share of weightings passing both of the original falsification tests is $0.0\%$ on dev and $41.2\%$ on test. The results do not support a ranking of composites.
\end{enumerate}
A statement from the earlier version survives: evidence-support measures that compare a response with the supplied evidence cannot detect corrupted evidence.
It was an argument then; Section~\ref{sec:results} makes it an experiment.
The full list of corrected claims is in the repository (\texttt{ERRATUM.md}).

\section{Discussion}

\paragraph{What the results do and do not say.}
The tests do not show that overlap, entailment or judge scores are useless.
They measure faithfulness to the supplied context, and they do that job: a response that invents a number is flagged.
The finding is about what that property cannot cover.
When the failure is in the evidence, any score defined relative to the evidence reports ``faithful'' and is wrong about reliability.
A composite that adds such a score to accuracy and calibration moves with corruption only through the two terms that depend on the gold answer, and so is no better than accuracy at a task where accuracy is unavailable.

\paragraph{What gets past the limit.}
Only checks that bring in something outside the response and the evidence it was given.
The cheapest such signal, redundancy inside the evidence, is weak and fails when numbers are moved rather than changed.
A trusted copy of the numbers handles changed numbers but not moved ones, which is the more realistic failure of a mis-parsed table.
The check that handled every error we tried compares (row, column, value) cells with a trusted structured source, such as the XBRL record behind a filing. It needs that source to exist and the evidence to carry table structure, which limits it to part of the items (a third on FinQA).
Where such a source exists the practical conclusion is simple: verify the evidence against it before generation, and do not read a faithfulness score as a reliability estimate.

\paragraph{Implications for practice and for evaluation.}
(i) A reliability pipeline needs a signal about the evidence, not only about the response.
(ii) A verifier that compares numbers is not enough; compare structure.
(iii) Validate an evidence-quality manipulation by its effect on the model before using it to evaluate a metric. The distractor sentences we first used would have been reported as a metric failure.

\section{Conclusion}

Reference-free reliability scores for numerical RAG cannot see wrong evidence, and a controlled test on two datasets and four generators from two model families confirms it: accuracy collapses while every measure that sees only the response and the supplied evidence stays near chance (AUROC at most 0.60).
Realistic errors that merely misplace genuine numbers cost as much accuracy and escape even a check of the numbers against the clean evidence; only a cell-level check against a trusted structured copy catches them.
The same response-side measures are weak and model-dependent at flagging wrong answers on clean evidence and real retrieval misses.
Practitioners who need reliability estimates should treat such scores as faithfulness checks and add a verification of the evidence against a trusted structured source.

\section*{Limitations}

The generators are open models of 1.5B to 14B parameters, run without fine-tuning; larger or finance-tuned models may behave differently, and the 14B model was run on 100 items.
Accuracy is low for the smaller models (about 0.2 on FinQA for the 3B model; 0.47 to 0.62 for the 7B and larger models), which limits the power of the correct-versus-incorrect comparisons; we report intervals instead of claiming equality.
The scaled-number corruption is synthetic and numeric. The realistic errors (swapped periods and rows) are closer to real parsing failures but are still injected, were run for one generator (Qwen2.5-7B) and only on table rows, and on FinQA affect only the quarter to third of items whose retrieved evidence has an eligible table row.
The cell-level check is, by construction, an upper bound: it assumes a trusted structured copy and that the evidence carries table structure.
The retrieval-miss analysis defines a miss by the absence of gold numbers in the retrieved text, which can disagree with whether the answer was recoverable; the TAT-QA gold evidence is the operands of the annotated derivation because the dataset has no cell-level evidence labels.
The entailment measure uses one small NLI model; the lexical baseline is a heuristic; the faithfulness judge is a single 4-bit 8B model with one prompt, run only on the clean and fully edited evidence, and a stronger judge or a different prompt may do better. It is a RAGAs-\emph{style} score, not the RAGAs implementation.
The 7B, 8B and 14B models are 4-bit builds run through a different inference engine than the float32 models, so differences between sizes are confounded with precision.
We do not report a GPU 4-bit run of the smaller models.
Item counts (100 to 200 per run) are modest; pooling over runs narrows the intervals but mixes models.

\section*{Reproducibility}
Code, prompts, per-response outputs for every run, and the GitHub Actions workflow that regenerates them on free CPU runners are at \url{https://github.com/aaronsunny007/Financial-NRF}. Tables and figures in this paper are produced from the result files by \texttt{experiments/paper\_analysis.py} and \texttt{experiments/paper\_figures.py}.

\section*{Acknowledgements}
The experiments, analysis code and a first draft of this text were produced with the assistance of Claude (Anthropic), working in an agentic coding environment. The author directed the study, reviewed the results and the text, and is responsible for all claims. No external funding supported this work.
